\documentclass[11pt]{article}
\usepackage[a4paper,margin=1in]{geometry}
\usepackage{amsmath,amssymb,amsthm}
\usepackage{booktabs}
\usepackage{graphicx}
\usepackage{xcolor}
\usepackage{hyperref}
\hypersetup{colorlinks=true,linkcolor=blue,urlcolor=blue,citecolor=blue}
\newcommand{\rt}[1]{\texttt{#1}}
\newcommand{\fn}[1]{\texttt{#1}}
\newcommand{\ifgraphic}[2]{\IfFileExists{#1}{\includegraphics[width=#2]{#1}}%
  {\fbox{\parbox{#2}{\centering\small figure \texttt{\detokenize{#1}} pending}}}}

\theoremstyle{plain}
\newtheorem{proposition}{Proposition}

\title{\textbf{Report R23 --- the stripes are walls, the flicker is a dark
qubit, and only rules 73 and 109 are doing many-body physics}\\
\large Krylov sector analysis of quantum cellular automata}
\author{Tier 1 analysis}
\date{2026-08-21}

\begin{document}
\maketitle

\begin{abstract}
The trajectory pictures for $W_{28}$, $W_{29}$, $W_{70}$, $W_{71}$, $W_{157}$
and $W_{199}$ all end in a stripe pattern whose gaps appear to breathe slightly,
which raises the fair suspicion that these six rules are trivial. They are ---
\emph{asymptotically}, and now for a stated reason rather than an impression.
Every terminal SCC of all six is a \textbf{subcube}: its states agree on a set of
frozen sites and range freely over the rest, so $|A|=2^{\#\mathrm{free}}$
(exhaustive, $N=6\ldots14$). Together with R17's completeness this says exactly
one thing --- the attractor is $k$ independent unbiased coins and nothing else.
The frozen part is a \textbf{wall tiling}: the attractor words are precisely the
tilings of the chain by \rt{10} and \rt{f10} (rules $199$, $71$; mirrored for
$157$, $29$) or by \rt{10}, \rt{1f0} plus a leading vacuum run (rules $28$,
$70$), as set equality for $N=6\ldots14$. That \emph{derives} R17's two fitted
constants: tiles of length $2$ and $3$ give the plastic number $\rho$, and the
all-wide tiling gives $2^{\lfloor(N+1)/3\rfloor}$. Every attractor of all eight
rules is moreover an exact decoherence-free subspace --- the jump probability on
it is $0$ identically --- so the channel restricted to an attractor is a
\emph{unitary}. For the six that unitary is a product of single-qubit Hadamards
on the free sites: each runs $H^t$ and oscillates with \textbf{period $2$}. The
stripes are the walls; the breathing gap is one dark qubit per wide tile; the
half-chain entanglement is identically zero because the attractor subspace is a
tensor product across every cut (membership-matrix rank $1$). The period-$2$
signal is single-shot: averaging over trajectories drops it from $\sim0.45$ to
$\sim0.05$, so it is invisible in an averaged plot. Rules $29$ and $71$ --- and
only they --- have weak components holding several terminal SCCs
($n_\mathrm{wcc}=n_\mathrm{rec}$ exactly for the other six). The selection is a
set of independent \textbf{fair coins resolved in one period}: absorption
probabilities are uniform on $2^k$ attractors with $k\in\{1,2\}$, never anything
else. It is multistability with noise-induced selection, not a bifurcation.
Finally $W_{73}$ and $W_{109}$ are a different animal and the answer to
``are they more interesting'' is yes: their attractors are \emph{not} subcubes
but the hard-core constrained spaces --- no two adjacent $1$s for $W_{73}$
($\dim=F(N+2)$), no two adjacent $0$s for $W_{109}$ ($\dim=F(N)$). The
dissipation prepares an exponentially large decoherence-free subspace carrying a
kinetically constrained unitary, and the state inside it is entangled.
\end{abstract}

\tableofcontents

\section{The question}

Rules $28$, $29$, $70$, $71$, $157$, $199$ are R17's plastic class. In the
trajectory simulations they all relax to a stripe pattern with what looks like a
slow oscillation of the space between stripes. Three things were asked: whether
that is really what happens, how it relates to the terminal-SCC digraph, and what
precisely goes on in $29$ and $71$, whose weak components do not have unique
terminal SCCs. A fourth question asked what $73$ and $109$ look like and whether
they are more interesting. All four are answered below; the last one has the
largest answer.

Everything here is at \rt{obc0} with $V=H$, and every statement labelled
\emph{exhaustive} is a full enumeration of $2^N$, not a sample.

\section{The stripes are walls}

\subsection{Every attractor is a subcube}

Write an attractor $A$ as a word over $\{0,1,\rt{f}\}$: site $i$ gets its common
value if every state of $A$ agrees there, and \rt{f} if it does not.

\begin{proposition}
For all six rules and $N=6,\dots,14$, every terminal SCC satisfies
$|A| = 2^{\#\rt{f}}$ --- $A$ is the full subcube on its free sites.
\end{proposition}

\noindent This is the missing half of R17. R17 showed each terminal SCC is a
complete digraph $K_n$; ``$K_n$ on a subcube'' is precisely the statement that
the free bits are redrawn independently and uniformly every period. The graph is
the product of $k$ unbiased coins, which is why its second eigenvalue is exactly
$0$ and why $n$ is always a power of two.

\subsection{The frozen part is a tiling, and it fixes both constants}

A tile is a wall \rt{1} together with the empty sites that belong to it; a wide
tile carries one free site.

\begin{proposition}
The set of attractor words is \emph{exactly} the set of tilings of the chain by
\[
\begin{array}{ll}
W_{199},\,W_{71} & \rt{10},\ \rt{f10}\\
W_{157},\,W_{29} & \text{the mirror image of those}\\
W_{28},\,W_{70}  & \rt{10},\ \rt{1f0},\ \text{preceded by a run of vacuum tiles \rt{0}}
\end{array}
\]
with two edge conditions. The chain may end inside a tile, but the remnant must
still contain its wall --- a free site with no wall beside it is not dark, so
\rt{f10} may be cut to \rt{f1} and never to \rt{f}. And for $W_{28}$, $W_{70}$
the chain must end \emph{at} a wall: a trailing empty site with nothing after it
is refilled, so the final tile is always the truncated one, whereas the $199$
family may end on a full tile. Set equality --- no containment either way ---
for $N=6,\dots,14$, and it is these two conditions that make it equality rather
than containment.
\end{proposition}

\noindent Two consequences drop out without a fit.

\emph{The plastic number.} Tiles of length $2$ and $3$ give the counting
recursion $a(N)=a(N-2)+a(N-3)$, whose growth rate is $\rho=1.3247$, the root of
$x^3=x+1$. The counts are $5,7,9,12,16,21,28,37,49$ for the $199$ family and
$11,15,20,27,36,48,64,85,113$ for the $28$ family over $N=6\ldots14$, matching
R17's series term for term. The $28$ family is larger by exactly the leading
vacuum run, and it is the \emph{absence} of an \rt{E} that permits it: with an
\rt{E} at $(0,0)$ the empty region refills, so $157$, $199$, $29$, $71$ have no
vacuum tile. That is R17's ``the children carrying an \rt{E} keep a strict
subset'', now with the reason attached.

\emph{The base $2^{1/3}$.} The largest attractor is the all-wide-tile word, so
its dimension is $2^{\lfloor (N+1)/3\rfloor}$ exactly, giving
$b_\mathrm{att}=2^{1/3}$ as an identity rather than a fitted constant.

\section{Every attractor is decoherence-free}

\begin{proposition}
For all eight rules, on every terminal SCC tested the jump probability is $0$ to
machine zero at every site of every period, the no-jump Kraus operator removes no
norm ($\le 3\times10^{-16}$), and the support never leaves the SCC. The channel
restricted to a terminal SCC is therefore a \emph{unitary}.
\end{proposition}

\noindent This is the sharp form of R17's ``the reset never fires on its own
attractor''. The dissipation is spent entirely on the transient; what it leaves
behind is a decoherence-free subspace. The interesting question then becomes
which unitary, and the six rules and the two rules answer it differently.

\section{The flicker: one dark qubit per wide tile}

For the six, the free sites are isolated --- each sits between a wall and an
empty site --- so the surviving unitary is a product of single-qubit Hadamards,
one per wide tile. Each free qubit runs $H^t$, and since $H^2=\mathbb{1}$ it
cycles $\lvert0\rangle\to\lvert+\rangle\to\lvert0\rangle$: the occupation at that
site alternates between a definite value and $1/2$ with \textbf{period $2$},
forever, exactly. Measured deep in the attractor,
$\max_t\lvert n(t+2)-n(t)\rvert = 0$ to machine precision at every $N$ tested and
in both brickwork conventions (engine even-first, \fn{pyqca} odd-first).

So the trajectory picture decomposes cleanly:

\begin{center}
\begin{tabular}{ll}
\toprule
what you see & what it is \\
\midrule
bright stripes, fixed & the walls of the tiling \\
dark gaps of width $1$ or $2$ & \rt{10} and \rt{f10} tiles \\
the gap that breathes with period $2$ & the dark qubit of a wide tile \\
which stripe pattern you get & which tiling the noise selected \\
\bottomrule
\end{tabular}
\end{center}

\subsection{The entanglement is exactly zero, not just small}

Being a subcube, the attractor subspace is a tensor product of single-site
spaces, hence a tensor product across \emph{every} cut. Checked directly at the
middle cut: the membership matrix of left-half against right-half configurations
has rank $1$ for all $41$/$77$ (rule $28$) and $20$/$36$ (rule $199$)
non-trivial attractors at $N=11,13$, and likewise for $70$, $157$, $29$, $71$. Hence \emph{every} state in the attractor is a product state, and the
measured half-chain entropy of the trajectory is $0$ identically --- not
statistically zero, deterministically zero, at every $N$ and on every trajectory.

\subsection{What survives averaging, and what does not}

Two different things happen to the two features, and they should not be lumped
together.

The \emph{oscillation} is destroyed. The period-$2$ amplitude of a single
trajectory is $0.42$--$0.50$, essentially its maximum of $1/2$; averaged over
$200$--$400$ trajectories it collapses to $0.03$--$0.07$. Different trajectories
land on different tilings and reach them at different times, so the phases add
incoherently.

The \emph{stripes} are damped but not erased. A single trajectory has full
contrast, $\max_i n_i - \min_i n_i = 1.000$: the frozen sites sit at exactly $0$
or $1$. Averaged over $300$ trajectories at $N=15$ the contrast falls to $0.48$
($W_{28}$) and $0.31$ ($W_{199}$), and what survives is not any one tiling but a
staggered residue --- a period-$2$ modulation in \emph{space} about $n\approx0.5$
with the boundary sites strongly pinned ($0.81$ at the right edge for $W_{28}$).
That residue is a consequence of the grammar: walls are separated by $2$ or $3$
sites, so across the tiling ensemble they still favour one sublattice, and the
open boundary fixes the first wall. For $73$ and $109$ the same average is much
flatter, contrast $0.12$ and $0.16$.

\textbf{Prediction for the notebook:} the cells plotting one trajectory should
show crisp full-contrast stripes with a visible flicker; the cell taking
\fn{np.mean} should show a much weaker staggered pattern about $n\approx1/2$ with
bright edges, and no time dependence at all. That is a check on this report, not
a result of it.

\begin{figure}[htbp]
\centering
\ifgraphic{../../figures/fig_stripes_obc0.pdf}{\textwidth}
\caption{\textbf{Top row:} one quantum trajectory each at $N=23$, the size used
in \fn{pyqca/mixed/trajectories.ipynb} (odd layer first, fixed-zero boundary).
$W_{28}$ and $W_{199}$ lock into a wall tiling; $W_{73}$ and $W_{109}$ never do.
\textbf{Middle row, left to right:} the same rule averaged over $300$
trajectories, where the stripes wash out because each run picks a different
tiling; the occupation of one free site, alternating with period $2$ against a
site of $W_{73}$ that does not; the period-$2$ amplitude single-shot against
ensemble-averaged; and the half-chain entanglement inside the attractor, which is
identically zero for the six and finite for $73$ and $109$.
\textbf{Bottom row:} the number of attractors against $N$ with $\rho^N$ and
$\psi^N$ for reference; the largest attractor against $2^{\lfloor(N+1)/3\rfloor}$
and $F(N{+}2)$; the fraction of attractors that are subcubes --- exactly $1$ for
the six, well below it for $73$ and $109$; and the number of enclosures per
strong-symmetry block, which is $1$ for every rule except $29$ and $71$.}
\label{fig:stripes}
\end{figure}

\section{Rules 29 and 71: several enclosures in one block}

\subsection{Only these two split}

Counting weak components against terminal SCCs at $N=9,\dots,13$:

\begin{center}
\begin{tabular}{lccccc}
\toprule
rule & $28$, $70$ & $157$, $199$ & $73$, $109$ & $29$ & $71$ \\
\midrule
$n_\mathrm{wcc}$ vs $n_\mathrm{rec}$ & equal & equal & equal &
  $7$ vs $16$ ($N{=}10$) & $12$ vs $16$ ($N{=}10$) \\
\bottomrule
\end{tabular}
\end{center}

\noindent Six of the eight have an exact bijection: one attractor per weak
component. The gap is a statement about \emph{which rules} can split, not about
every size: it is absent below $N=6$ for $W_{29}$ (which has it at $N=4$, loses
it at $N=5$, and keeps it from $N=6$ on) and below $N=7$ for $W_{71}$. Over
$N=7,\dots,14$ both show it at every size. Cross-checked against the Tier-1
session's independent enumeration at $N=4,5$, which agrees row for row. In the R22 dictionary a weak component is a strong-symmetry block and
a terminal SCC is an enclosure, so those six have a \emph{unique} steady state
inside every block. Rules $29$ and $71$ do not: a block can carry up to six
enclosures, and the asymptotic state depends on the trajectory.

\subsection{What happens is a fair coin, resolved in one period}

\begin{proposition}
For $W_{29}$ and $W_{71}$ at $N=8,\dots,12$: every state whose fate is undecided
has \emph{all} of its successors already committed, so the branch resolves in
exactly \textbf{one period}; and the absorption distribution is uniform on $2^k$
attractors with $k\in\{1,2\}$ --- probability vectors are $(\tfrac12,\tfrac12)$
($560$ states at $N=12$) or $(\tfrac14,\tfrac14,\tfrac14,\tfrac14)$ ($8$ states),
and nothing else ever occurs.
\end{proposition}

\noindent The shared-route states are $1.6\%$--$14\%$ of the space
($4,24,32,128,192$ for $W_{71}$ at $N=8\ldots12$) and are all transient.

\subsection{What the coin chooses, and why it needs both resets}

Take $W_{29}$ at $N=9$ and the branch point \rt{110001000} (site $0$ leftmost).
In one period it reaches eight states with probability $1/8$ each, four in each
of two attractors. Tracing the layers: on the even layer site $2$ sees
$(m,n)=(1,0)$ and gets a Hadamard --- that is the coin --- and site $8$ sees
$(0,0)$ and is excited by the \rt{E}. On the odd layer the outcome at site $2$
routes everything else:
\begin{itemize}
\item site $2$ measures $1$: site $3$ then sees $(1,0)$ and is Hadamarded into
the free site of a wide tile, while site $1$ sees $(1,1)$ and is \emph{reset by
the \rt{D}}. The wall lands at site $2$, giving the tiling $(2,3,3,1)$.
\item site $2$ measures $0$: site $3$ sees $(0,0)$ and is \emph{excited by the
\rt{E}} into a new wall, while site $1$ is Hadamarded into the free site. The
wall lands at site $3$, giving the tiling $(3,2,3,1)$.
\end{itemize}
So the coin exchanges a short tile with a wide one, and it needs \emph{both}
resets: the
\rt{E} writes the new wall, the \rt{D} erases the old one. That is exactly what
$29=\rt{EIVD}$ and $71=\rt{EVID}$ have and what $28=\rt{IIVD}$, $70=\rt{IVID}$,
$157=\rt{EIVI}$, $199=\rt{EVII}$ do not --- which is why only these two split.

\subsection{What survives as a quantum number}

The block is not label-free. Across all $20$ multi-attractor blocks of $W_{29}$
and all $16$ of $W_{71}$ at $N=8,\dots,12$, the \textbf{number of walls is
constant} --- that is the strong symmetry the extra reset leaves intact. Their
\emph{arrangement} is not: the coin permutes tiles. For $W_{71}$ the multiset of
tile lengths is conserved as well ($16/16$), so the branch is a pure
rearrangement, e.g.\ a block whose three attractors are the tilings
$(3,2,2,1)$, $(2,3,2,1)$, $(2,2,3,1)$. For $W_{29}$ the multiset survives in only
$8/20$ blocks, because its truncated boundary tile can trade a unit --- one block
at $N=10$ holds $(2,2,2,2,2)$ alongside $(2,2,3,2,1)$, same five walls,
different widths. The two rules are a reflection pair and \rt{obc0} with a fixed
layer order is not reflection-covariant for dissipative rules, so an asymmetry at
the right edge is expected rather than surprising.

\subsection{Is it a bifurcation? No}

A bifurcation is a qualitative change of the attractor set as a parameter is
varied. Nothing is varied here: there is no coupling to tune, the branching ratio
is exactly $1/2$ at every $N$, and the attractor set is the same tiling family as
for $157$ and $199$ --- rules $29$ and $71$ have \emph{identical} attractor words
to $157$ and $199$ respectively; the extra \rt{D} changes only which basins are
joined. The correct name is multistability with noise-induced selection: several
enclosures inside one strong-symmetry block, chosen by the measurement record.

Nor is the branch coherent. Starting the unmonitored trajectory \emph{at} the
branch point, $40$ runs all end supported in a single terminal SCC with
$S_\mathrm{mid}=0$ ($19$ one way, $21$ the other) --- the \rt{D} acts as a
which-path detector, so the ensemble is a classical $\tfrac12/\tfrac12$ mixture
of two stripe patterns, not a superposition of them.

\subsection{Coherence between attractors survives too}

\S3 was stated one attractor at a time, which understates it. The no-jump Kraus
operator is a single global operator serving every enclosure at once, so by
linearity it annihilates the jump branch on the whole recurrent space, not merely
on each piece. Checked directly: prepare
$(\lvert a\rangle+\lvert b\rangle)/\sqrt2$ with $a$ and $b$ drawn from
\emph{different} terminal SCCs, and over $20$ periods the jump probability stays
$0$ exactly, the support stays in $A\cup B$, and the weight stays at
$0.5000/0.5000$ --- for all six rules tested. A superposition of two different
stripe patterns is preserved; the machine's own resets never decohere it.

Assembling the induced map on the whole recurrent space makes the period-$2$
result algebraic rather than numerical. For the six, that map has \emph{exactly
two} distinct eigenvalues and satisfies $U^2=\mathbb{1}$ to
$2.2\times10^{-16}$ --- e.g.\ $D=9$ with multiplicities $(6,3)$ for $W_{28}$ at
$N=4$, $D=11$ with $(6,5)$ for $W_{199}$ at $N=6$. Period $2$ is then forced, not
measured. For $W_{73}$ at $N=5$ the same map has $12$ distinct eigenvalues with
multiplicities $(8,4,2,2,1^8)$ --- stable under clustering tolerances from
$10^{-10}$ to $10^{-2}$ --- and $\lVert U^2-\mathbb{1}\rVert=0.5$, so no period,
as \S6.4 found.

This also settles which of these rules could hold a code. It cannot be any of
them: a subspace spanned by computational basis states cannot detect a
single-site $Z$, because the $W^\dagger Z_iW$ are diagonal with $\pm1$ entries
given by the $i$-th bit and distinct codewords are distinct bit strings, so every
joint eigenspace is one-dimensional. All eight attractors are basis-spanned ---
subcubes for the six, constrained subsets for $73$ and $109$ --- so all have
distance $1$ and no subcode helps. The honest description is a noiseless
subsystem against \emph{this machine's own resets} and nothing else. (Argument
due to the Tier-1 session, R24.)

\section{Rules 73 and 109 are doing something else}

\subsection{The attractors are constrained spaces, not subcubes}

Only $13/55$ and $21/68$ of the non-trivial attractors of $W_{73}$ and $W_{109}$
at $N=11$ are subcubes, against $100\%$ for the six. The largest one is the
diagnostic:

\begin{center}
\begin{tabular}{lll}
\toprule
rule & largest attractor & dimension \\
\midrule
$W_{73}$  & all bitstrings with \textbf{no two adjacent $1$s} & $F(N+2)$ \\
$W_{109}$ & all bitstrings with \textbf{no two adjacent $0$s} & $F(N)$ \\
\bottomrule
\end{tabular}
\end{center}

\noindent Verified for $N=7,\dots,13$: $W_{73}$ gives $34,55,89,144,233,377,610$
and $W_{109}$ gives $13,21,34,55,89,144,233$. This is the hard-core / blockade
constrained space, and it is where $b_\mathrm{att}=\varphi$ comes from. The
number of attractors grows at the supergolden ratio $\psi=1.4656$
($13,19,28,41,60,88,129$ for $W_{73}$).

Combined with \S3, the physical statement is: \textbf{the dissipation prepares an
exponentially large decoherence-free subspace which is a kinetically constrained
Hilbert space, and then acts as a unitary constrained QCA inside it.} That is the
setting of PXP-type models, and it is not something the six rules do.

\subsection{The induced unitary is PXP with the spin flip replaced by a Hadamard}

\S3 says the channel restricted to an attractor is \emph{some} unitary. For these
two it is an identifiable one.

\begin{proposition}
On the largest attractor of $W_{73}$ and of $W_{109}$, at $N=9$ and $11$, the
induced one-period map equals, to machine precision
($\le1.7\times10^{-16}$ entrywise), the brickwork product over the even then the
odd sublattice of the blockade-controlled Hadamard
\[
  U_i \;=\; P^{a}_{i-1}\,H_i\,P^{a}_{i+1}
        \;+\; \bigl(\mathbb{1}-P^{a}_{i-1}P^{a}_{i+1}\bigr),
\]
with $a=0$ (fire only when both neighbours are \emph{empty}) for $W_{73}$ and
$a=1$ (both neighbours \emph{occupied}) for $W_{109}$; sites beyond the chain
count as $0$. The map is unitary on that subspace to
$\lVert AA^{\mathsf T}-\mathbb{1}\rVert\le4.4\times10^{-16}$.
\end{proposition}

\noindent This is the PXP structure exactly --- a projector-dressed single-site
operator on a hard-core constrained chain --- with the spin flip $X$ replaced by
the Hadamard $H$, and in Floquet (brickwork) rather than Hamiltonian form. The
constraint is preserved for the obvious reason: flipping a site whose neighbours
are both empty can never create an adjacent pair, so $P^0HP^0$ maps the
no-two-adjacent-$1$s space to itself, and $P^1HP^1$ does the same for
no-two-adjacent-$0$s.

The two rules are \emph{not} related by a particle-hole conjugation, which is
worth stating because the constraints look like mirror images. Conjugating by a
global $X$ would send $H\mapsto XHX$, and $XHX$ is a different reflection: the
$W_{109}$ map tested against a $P^1(XHX)P^1$ brickwork misses by $0.71$, while
the same test with $H$ is exact. Both rules carry the \emph{same} gate and
opposite blockade projectors, and the \rt{obc0} boundary pins $0$ rather than
$1$, which is why the dimensions come out at $F(N+2)$ and $F(N)$ rather than
agreeing.

\subsection{And they are entangled}

The membership matrix across the middle cut has rank $2$ --- the blockade
constraint is a bond-dimension-$2$ condition, so the attractor is \emph{not} a
product subspace, and states in it need not be product states. They are not: the
half-chain entropy of the trajectory state deep in the attractor is

\begin{center}
\begin{tabular}{lcccc}
\toprule
 & $N=9$ & $11$ & $13$ & $15$ \\
\midrule
$W_{73}$  & $0.624\pm0.121$ & $0.737\pm0.158$ & $0.773\pm0.178$ & $0.928\pm0.215$ \\
$W_{109}$ & $0.386\pm0.084$ & $0.616\pm0.117$ & $0.737\pm0.152$ & $0.751\pm0.166$ \\
\midrule
the six & $0$ & $0$ & $0$ & $0$ \\
\bottomrule
\end{tabular}
\end{center}

\noindent ($24$ trajectories each, averaged over an $8$-period window, errors are
standard errors of the mean; the six rules read $\le 9.3\times10^{-17}$, i.e.\
machine zero, at every entry.) The contrast with that row of exact zeros is the
result. The apparent \emph{growth} with $N$ is \textbf{not} resolved at this
sample size --- the rise from $N=9$ to $N=15$ is $1.2\sigma$ for $W_{73}$ and
$2.0\sigma$ for $W_{109}$ --- so we deliberately fit no scaling law to it. The
claim is that the entropy is nonzero and $O(1)$, not that it obeys a volume law.
A run with enough trajectories to separate a volume law from a logarithm is the
obvious next measurement and is not attempted here.

Their occupation profile also has no period: $\max_t\lvert n(t+p)-n(t)\rvert$
stays $O(1)$ for every $p\le8$, against machine zero at $p=2$ for the six.

\subsection{This is the structural counterpart of the MPO result}

The dissipative MPO sweep found an $N$-independent operator Schmidt rank for all
six --- $\chi_\mathrm{sat}=53$ ($28$, $70$), $18$ ($29$, $71$), $67$ ($199$,
$157$) --- and exponential growth for $73$ and $109$. The structure above is why.
The asymptotic map of the six is a sum over a \emph{regular} tiling language of
\emph{product} operators, and a regular language has a finite transfer matrix, so
the bond dimension cannot grow with $N$. For $73$ and $109$ the asymptotic map
carries an entangling unitary on a constrained space, and nothing bounds it. We
state this as the mechanism, not as a derivation: the particular values $18$,
$53$, $67$ are not accounted for here.

\section{Answering the question as asked}

\begin{enumerate}
\item \textbf{Is the stripe pattern real?} Yes, and it is a wall tiling. The
stripes are the walls, the gaps are tiles of width $2$ or $3$, and which tiling
you get is chosen by the noise from a set that grows like $\rho^N$.

\item \textbf{Is the oscillation real?} Yes, and it is exactly period $2$: one
dark qubit per wide tile running $H^t$, coherent and immune to the dissipation
because the attractor is decoherence-free. It is a single-shot effect and
averages away.

\item \textbf{How does it relate to the terminal-SCC digraph?} Directly. The
terminal SCC is a complete digraph on a subcube, and that is the graph-theoretic
shadow of ``$k$ independent dark qubits'': completeness $=$ the bits are
redrawn independently and uniformly; the subcube $=$ they are independent
\emph{sites}; the frozen complement $=$ the walls.

\item \textbf{Are the six trivial?} Asymptotically yes --- product state, zero
entanglement, memoryless, a decoupled period-$2$ flicker. What is not trivial is
the \emph{fragmentation}: an exponential number of decoherence-free subspaces
surviving dissipation, selected by the measurement record. The interest is in the
selection, not in the asymptotic dynamics.

\item \textbf{What happens in $29$ and $71$?} A weak component holds several
enclosures. From a shared state the fate is settled in one period by $k$
independent fair coins, uniform on $2^k$ outcomes, $k\in\{1,2\}$. The coin picks
the width of one tile and requires both an \rt{E} and a \rt{D}. Not a
bifurcation: no parameter, ratio pinned at $1/2$, same attractor set as $157$ and
$199$.

\item \textbf{Are $73$ and $109$ more interesting?} Decisively. Their attractors
are the hard-core constrained spaces of dimension $F(N+2)$ and $F(N)$ --- not
product subspaces, rank $2$ across a cut, hosting an entangled and aperiodic
constrained unitary. The induced unitary is
identified exactly: blockade-controlled Hadamards, PXP with $X\to H$ in Floquet
form. Dissipative preparation of an exponentially large decoherence-free
subspace carrying a PXP-type constrained dynamics is worth a paper; a set of
decoupled flickering qubits is not.
\end{enumerate}

\section{Reproduction}

\begin{verbatim}
python -m qca_fragmentation.scaling.stripes --rebuild
python -m pytest qca_fragmentation/tests/test_stripes.py -q
\end{verbatim}

\noindent All numbers come from \fn{analytics/stripes\_obc0.json}. The exact
statements (subcube, tiling grammar, absorption law, decoherence-free test) are
full enumerations of $2^N$; the entanglement and period-$2$ numbers are
state-vector quantum trajectories with the per-site two-Kraus split of
\fn{core.cycle.\_split\_site}, sampled.

\end{document}
